% ============================================================================
\section{Realidad problemática}

La elaboración de estadísticas macroeconómicas ocurre bajo restricciones prácticas de información y sucesivos procesos de actualización. En el caso del PBI mensual, el primer dato difundido se construye con registros preliminares que se amplían y concilian posteriormente. Por esta razón, las revisiones no son simples errores aleatorios independientes, sino ajustes que reflejan ampliaciones de cobertura, llegada tardía de encuestas y cambios metodológicos. Entender este mecanismo resulta esencial porque la toma de decisiones económicas se realiza con las cifras disponibles en cada momento, y no con los valores corregidos años después.

La literatura especializada advierte que evaluar modelos económicos usando la versión revisada más reciente introduce un sesgo retrospectivo, ya que presupone información que no existía cuando se tomó la decisión \parencite{croushore2011}. Por ello, estudiar revisiones exige reconstruir la secuencia cronológica exacta de publicaciones —conocidas como \emph{vintages}— conservando en cada fecha únicamente los datos ya divulgados.

En el contexto peruano, antes de afirmar si las revisiones del PBI tienen dirección sistemática o son predecibles, resultó imprescindible auditar directamente las fuentes históricas, sus fechas de difusión y la homogeneidad de las definiciones sectoriales. Sin comprobar que cada dato corresponde inequívocamente a su fecha real de publicación, cualquier comparación predictiva pierde validez empírica.

Un segundo desafío surge en el plano metodológico al combinar series temporales relativamente cortas con modelos altamente flexibles, como las arquitecturas cuánticas variacionales o de reservorio. Aunque estas herramientas producen transformaciones no lineales complejas, dicha flexibilidad no garantiza mejores pronósticos fuera de muestra: en muestras macroeconómicas mensuales limitadas puede generar sobreajuste, dispersión numérica o resultados inferiores frente a modelos clásicos simples y de menor costo. En consecuencia, este estudio limitó el contraste cuántico confirmatorio a una única arquitectura de reservorio cuántico (QRC), excluyendo variantes variacionales más costosas para mantener una comparación simétrica y controlada.

\subsection{Evidencia descriptiva del repositorio auditado}

La auditoría de los archivos históricos del Banco Central de Reserva del Perú (BCRP) compilados por \textcite{cruz2025} permitió verificar la integridad de la información tras depurar \(3\,131\) registros repetidos exactos presentes entre 1994 y mediados de 2024 y comprobar la ausencia total de filtraciones futuras (véase el Anexo~\ref{app:auditoria}). En total se reconstruyeron \(44\,307\) observaciones mensuales por rama económica para nueve industrias (agricultura, comercio, construcción, electricidad, pesca, PBI agregado, manufactura, minería y servicios), abarcando publicaciones continuas desde enero de 1994 hasta diciembre de 2025 (\(384\) meses) sobre períodos de referencia entre enero de 1992 y octubre de 2025 (\(406\) meses). Asimismo, el rezago habitual de publicación entre el mes medido y la primera cifra difundida se concentró sistemáticamente en dos meses (media de \(2{,}003\) meses, con julio de 2019 como único registro a tres meses).

Esta revisión exhaustiva permitió resolver tres problemas que distorsionaban las lecturas preliminares de la base. Primero, la convención inicial de restar la primera cifra menos la última tenía el signo contrario a la definición de revisión acumulada (Ecuación~\ref{eq:target}): una diferencia media negativa de \(-0{,}256\) puntos porcentuales indicaba en realidad que la primera estimación suele quedar \emph{por debajo} de la cifra corregida, es decir, una revisión promedio al alza de \(+0{,}256\) puntos porcentuales. Segundo, al excluir los \(21\) meses iniciales sin publicación oportuna (1992--1993), los últimos meses sin cierre y los \(24\) períodos afectados por los cambios de año base de julio de 2000 y marzo de 2014, quedaron \(350\) observaciones limpias a un horizonte continuo de once meses (media de \(+0{,}347\) puntos porcentuales y desviación estándar de \(0{,}768\)), frente a \(331\) observaciones a doce meses, donde se pierden \(18\) registros internos por el formato móvil de las Notas Semanales. Tercero, se constató que la cobertura por sectores no es uniforme: servicios comienza en julio de 1997, electricidad recién en mayo de 2003 y minería presenta un vacío de nueve publicaciones entre agosto de 2018 y abril de 2019.
